\section{INTRODUCCIÓN}


En la introducción se menciona claramente el para qué y el porqué del documento, se incluye el planteamiento del problema, el objetivo, preguntas de investigación, la justificación.

No utilice en el documento la primera persona en singular (yo realicé las encuestas) ni primera persona plural (realizamos las encuestas); utilice siempre la narración en tercera persona (se realizaron las encuestas, se publicaron resultados, se establecieron parámetros, etc.).\footnote{No utilice los pies de página para citas bibliográficas. Los pies de página se utilizan para complementar información del texto, procure que sean fragmentos cortos para no distraer o confundir al lector.}

No menos importante es la utilización de conectores que unen elementos de una oración, tener una buena variedad de estos enriquecen la estructura y redacción del texto. Algunos ejemplos:

\begin{tabular}{ll}
    Sin embargo         & En conclusión \\
    Puesto que          & En pocas palabras \\
    Por consiguiente    & A continuación \\
    Dado que            & Acto seguido \\
    Teniendo en cuenta  & Con motivo de \\
    Entonces            & A saber \\
    Simultáneamente     & De la misma forma \\
    Posiblemente        & En síntesis \\
    En efecto           & Así \\
    Ya que              & Para concluir \\
    Ahora bien          & Luego \\
    En cambio           & Resumiendo \\
    En cuanto a         & De igual manera \\
    El siguiente punto es& Al mismo tiempo \\
    Así pues            & Probablemente \\
    Recapitulando       & Indiscutiblemente 
\end{tabular}

La presente plantilla de \LaTeX, como cualquier documento de formato \textit{.tex} necesita un editor de \LaTeX\ ya sea instalado en su dispositivo o un editor online. Como todos los archivos se necesitan entre sí, todos los que aquí se encuentran debe ser subidos o leídos por su editor, además de la carpeta con las imágenes ya que algunos comandos usan la ruta de ella (del tipo: imgs/imagen.jpg)\footnote{Si tienes dudas, solicita asesoría en biblioteca, CRAI o centro de documentación respectivo}. Además, recomendamos tener especial cuidado con los archivos \textit{macros.tex} y \textit{ieeeudea.bst}, estos dos soportan toda la estructura general de la plantilla y los parámetros de las normas IEEE; cualquier modificación, por mínima que sea de ellos, puede generar severos problemas en la compilación.

Está plantilla de \LaTeX \ fue creada con divisiones en secciones, donde el primer nivel es la división \textit{$\setminus$section\{\}} , la de segundo nivel \textit{$\setminus$subsection\{\}}, \textit{$\setminus$subsubsection\{\}} y \textit{$\setminus$paragraph\{\}} los dos siguientes niveles; asegúrese usar esta división para que se puedan cargar cada partición en la tabla de contenido de forma automática; si escribe los diferentes títulos sin estos comandos, deberá cosntruir la tabla de contenido. Recuerde que el \LaTeX\  también tiene la herramienta de divisiones con *, las cuales aparecen en el documento con su respectivo título y no aparecen en la tabla de contenido.